\documentclass[10pt,a4paper]{amsart}
\usepackage[margin=19mm]{geometry}
\usepackage{amsmath,amssymb,mathtools}
\usepackage[scr=rsfs,cal=euler]{mathalfa}
\usepackage{xfrac}
\usepackage[unicode,hidelinks,bookmarks=false]{hyperref}
\newcommand{\ie}{i.e.\ }
\newcommand{\cf}{cf.\ }
\newcommand{\resp}{respectively\ }
\newcommand{\Subsection}{\subsection}
\providecommand{\nobreakdash}{\nobreak\hbox{-}}
\setlength{\emergencystretch}{2em}
\multlinegap0pt
\usepackage{amssymb}
\usepackage{xcolor}
\usepackage[shortlabels]{enumitem}
\usepackage{bm}
\newtheorem{theorem}{Theorem}
\newtheorem{lemma}{Lemma}
\newcommand{\A}{\mathcal{A}}
\newcommand{\B}{\mathcal{B}}
\newcommand{\N}{\mathbb{N}}
\newcommand{\colour}{\mathtt{color}}
\newcommand{\ds}{\displaystyle}
\newcommand{\lan}{\mathcal{L}}
\title{New examples of words for which binomial complexities and subword complexity coincide}
\author{L\'eo Vivion}
\date{Source: arXiv:2509.11172v2 (CC BY 4.0)}
\begin{document}
\maketitle

\noindent\emph{Source and licence.} New examples of words for which binomial complexities and subword complexity coincide. Version arXiv:2509.11172v2 (CC BY 4.0), distributed by arXiv under the Creative Commons Attribution 4.0 International licence, http://creativecommons.org/licenses/by/4.0/, as recorded in the arXiv licence metadata for this exact version and shown on the abstract page https://arxiv.org/abs/2509.11172v2. Full licence notice: \texttt{licenses/arxiv-2509.11172-CC-BY-4.0.txt}.\par\smallskip

\noindent\textit{Editorial scope.} Counted unit: the article's lemma lemma:stability (source0.tex lines 393--395). The article's complete original proof of it is delivered with it (source0.tex lines 619--669, one \texttt{proof} environment of 51 lines, which the article places away from the statement). No mathematics is written, repaired or re-derived: the statement and the proof are the article's own lines, in the article's own notation, and no retained line is substituted or re-flowed.  Retained uncounted as the source-local context the proof needs: lines 148--154; lines 340--342; lines 486--486; lines 502--504.  Nothing was omitted from any retained span, so the package carries no omission record.  No retained line contains a citation, so the wrapper carries no bibliography and no bibliography entry.  No retained line contains a figure environment, an image-inclusion command or drawing code, so no image is delivered and the package declares no asset dependency.  Editorial formatting only, in addition: an AMS \texttt{amsart} preamble that restates the article's own declarations and macros used by the retained lines, namely the environments \texttt{theorem}, \texttt{lemma} on separate counters, together with the standard packages those lines need.  The retained environments print in this document as Theorem 1, Lemma 1 in order of appearance, whereas the article prints the same items with the numbering of its own full document.  The complete native source of the article as distributed in the arXiv e-print for this exact version is bundled unchanged as \texttt{sources/185019/source0.tex}.
\begin{theorem}\label{th:main}
Let $d\geq 2$.\\
\emph{(i)} If $w$ is a $1$-balanced $d$-ary word, then its $2$-binomial complexity is equal to its subword complexity.\\
\emph{(ii)} If $w$ is a word with subword complexity $n\in\N_{>0}\mapsto n+(d-1)$, then its $2$-binomial complexity is equal to its subword complexity.\\
\emph{(iii)} If $w$ is a hypercubic billiard word in dimension $d$, then its $2$-binomial complexity is equal to its subword complexity.\\
\emph{(iv)} If $w$ is a $d$-ary Sturmian coloring, then its $2$-binomial complexity is equal to its subword complexity. 
\end{theorem}

\begin{lemma}\label{lemma:projections_kbin}
Let $d\geq 2$, $k\in\N_{>0}$ and $w\in\{1,\ldots,d\}^\N$. If, for every pair of distinct letters $i,j\in\{1,\ldots,d\}$, $b_{\pi_{i,j}(w)}^k=p_{\pi_{i,j}(w)}$, then the $k$-binomial complexity of $w$ coincides with its subword complexity.
\end{lemma}

\section{Proof of Lemma~\ref{lemma:projections_kbin} and of Theorem~\ref{th:main}, assertions (i)--(iii)}\label{sec:binary_projections}

\begin{equation}\label{eq:projections}
    \binom{\pi_{i,j}(u)}{x}=\binom{u}{x} \hspace{0.3cm}\text{ and }\hspace{0.3cm} \binom{\pi_{i,j}(v)}{x}=\binom{v}{x}.
\end{equation}

\begin{lemma}\label{lemma:stability}
Let $w_0\in\A^\N$ and $w_1\in\B^\N$ be two infinite words written over disjoint alphabets. If there exists $k\in\N_{>0}$ such that $b_{w_0}^k=p_{w_0}$ and $b_{w_1}^k=p_{w_1}$, then, for every letter $a\in\A$, the $k$-binomial complexity of $\colour(w_0,a,w_1)$ coincides with its subword complexity.
\end{lemma}

\begin{proof}[Proof of Lemma~\ref{lemma:stability}]
Let $w_0\in\A^\N$ and $w_1\in\B^\N$ be two infinite words defined over disjoint finite alphabets. Let $a\in\A$ and $w:=\colour(w_0,a,w_1)$ be the coloring of the letter $a$ in $w_0$ by $w_1$. We assume that $b_{w_0}^k=p_{w_0}$ and $b_{w_1}^k=p_{w_1}$ for some $k\in\N_{>0}$ and aim to prove that $b_w^k=p_w$. To this end, consider two factors $u,v\in\lan(w)$ such that $u\sim_k v$. Our goal is to show that $u=v$.

We denote by $\sigma$ the substitution that maps $w$ back to $w_0$:
\[
    \begin{array}{rccl}
        \sigma:&\A\sqcup\B\setminus\{a\}&\longrightarrow&\A\\
        & i & \longmapsto & \begin{cases} i &\text{if $i\in\A$,} \\ a &\text{if $i\in\B$.}\end{cases}
    \end{array}
\]
Clearly, $\sigma(w)=w_0$, $\sigma(u),\sigma(v)\in\lan(w_0)$, $\pi_\B(w)=w_1$ and $\pi_\B(u),\pi_\B(v)\in\lan(w_1)$. Moreover,
\[
    u=\colour(\sigma(u),a,\pi_{\B}(u)) \hspace{0.5cm} \text{ and }\hspace{0.5cm} v=\colour(\sigma(v),a,\pi_\B(v)).
\]
Therefore, to prove that $u=v$ it is suficient to prove that $\pi_{\B}(u)=\pi_\B(v)$ and $\sigma(u)=\sigma(v)$ (it is even equivalent). Thanks to a reasoning similar to that used in the proof of Lemma~\ref{lemma:projections_kbin}, we have $\pi_\B(u)\sim_k\pi_\B(v)$ (see equation \eqref{eq:projections} in Section~\ref{sec:binary_projections}), and then $\pi_\B(u)=\pi_\B(v)$. We now prove that $\sigma(u)=\sigma(v)$. To begin with, we claim that for any $x\in\A^*$
\begin{equation}\label{eq:coloring}
    \binom{\sigma(u)}{x}=\sum\limits_{\substack{y\in(\A\sqcup\B\setminus\{a\})^{|x|} \\ \sigma(y)=x}} \binom{u}{y}.
\end{equation}
We temporarily postpone the proof of this claim. Since $u\sim_k v$, we have $\binom{u}{y}=\binom{v}{y}$ for every $y\in(\A\sqcup\B\setminus\{a\})^{\leq k}$. Combined with \eqref{eq:coloring}, this yields $\binom{\sigma(u)}{x}=\binom{\sigma(v)}{x}$ for every $x\in\A^{\leq k}$, \emph{i.e.}, $\sigma(u)\sim_k\sigma(v)$. Finally, since both $\sigma(u)$ and $\sigma(v)$ belong to $\lan(w_0)$, and since the $k$-binomial complexity of $w_0$ coincides with its subword complexity, it follows that $\sigma(u)=\sigma(v)$.

\medskip

To conclude, it only remains to prove relation \eqref{eq:coloring}. Writing $u=u_1u_2\ldots u_p$, where each $u_i$ is a letter in $\A\sqcup\B\setminus\{a\}$, we have $\sigma(u)=\sigma(u_1)\sigma(u_2)\ldots\sigma(u_p)$ where each $\sigma(u_i)\in\A$ is also letter. Therefore, by definition of binomial coefficients of words, we have
\[
    \begin{array}{c}
        \ds\binom{\sigma(u)}{x}:=\#\big\{(i_1,i_2,\ldots,i_{|x|}) \;\big |\; 1\leq i_1<i_2<\ldots<i_{|x|}\leq p \text{ and } \sigma(u_{i_1})\sigma(u_{i_2})\ldots\sigma(u_{i_{|x|}})=x\big\},\\
        \\
        \ds\binom{u}{y}:=\#\big\{(j_1,j_2,\ldots,j_{|y|}) \;\big |\; 1\leq j_1<j_2<\ldots<j_{|y|}\leq p \text{ and } u_{j_1}u_{j_2}\ldots u_{j_{|y|}}=y\big\}.
    \end{array}
\]
Hence, the identity \eqref{eq:coloring} is an immediate consequence of the set equality
\begin{equation}\label{eq:coloringb}
    E(u,x)=\bigsqcup\limits_{\substack{y\in(\A\sqcup\B\setminus\{a\})^{|x|} \\ \sigma(y)=x}} F(u,y)
\end{equation}
where
\[
    \begin{array}{l}
        E(u,x):=\big\{(i_1,i_2,\ldots,i_{|x|}) \;\big |\; 1\leq i_1<i_2<\ldots<i_{|x|}\leq p \text{ and } \sigma(u_{i_1})\sigma(u_{i_2})\ldots\sigma(u_{i_{|x|}})=x\big\},\\
        F(u,y):=\big\{(j_1,j_2,\ldots,j_{|y|}) \;\big |\; 1\leq j_1<j_2<\ldots<j_{|y|}\leq p \text{ and } u_{j_1}u_{j_2}\ldots u_{j_{|y|}}=y\big\}.
    \end{array}
\]
Let us prove that \eqref{eq:coloringb} holds. First, it is clear that the sets $(F(u,y))_y$ are disjoint. Indeed, if there exist $y,z\in (\A\sqcup\B\setminus\{a\})^{|x|}$ such that $F(u,y)\cap F(u,z)\neq\emptyset$, then there exists $(j_1,\ldots,j_{|x|})$ such that $y=u_{j_1}\ldots u_{j_{|x|}}=z$. Secondly, if $(i_1,\ldots,i_{|x|})\in E(u,x)$, then $\sigma(u_{i_1})\ldots\sigma(u_{i_{|x|}})=x$. Consequently, $y:=u_{i_1}\ldots u_{i_{|x|}}$ belongs to $(\A\sqcup\B\setminus\{a\})^{|x|}$ and satisfies $\sigma(y)=x$. In other words
\[
    (i_1,\ldots,i_{|x|})\in F(u,y)\subseteq \bigsqcup\limits_{\substack{z\in(\A\sqcup\B\setminus\{a\})^{|x|} \\ \sigma(z)=x}} F(u,z).
\]
Finally, if $(j_1,\ldots,j_{|x|})\in F(u,y)$ for some $y\in(\A\sqcup\B\setminus\{a\})^{|x|}$ satisfying $\sigma(y)=x$, then
\[
    x=\sigma(y)=\sigma(u_{j_1}\ldots u_{j_{|x|}})=\sigma(u_{j_1})\ldots\sigma(u_{j_{|x|}}),
\]
which means that $(j_1,\ldots,j_{|x|})$ belongs to $E(u,x)$. This completes the proof.
\end{proof}

\end{document}
